% GENERATED FILE. Edit canonical lessons, then run npm run generate:books.
% canonical-source-hash: fnv1a64:3acd7a70eb96bfaa
% canonical-lessons: KA-C526-bandhu, KA-C526-tavaru, KA-C526-avali, KA-C526-marimaga, KA-C526-kiriya

\chapter{Kin}
\label{ch:ka-kin}
\hlchaptermodality{526}

\begin{chapteropening}
\textbf{By the end of this chapter:} \emph{I can say kin, a woman's parental home, twins, a great-grandson and younger.}

Complete the last lesson of chapter 526: I can say kin, a woman's parental home, twins, a great-grandson and younger.
\end{chapteropening}

\section[\texorpdfstring{\kn{ಬಂಧು} (bandhu)}{ಬಂಧು (bandhu)}]{\kn{ಬಂಧು} (bandhu) — kin; a relative}

\label{lesson:KA-C526-bandhu}



\begin{quote}
Before the new one: say the Kannada for a cousin, then the Kannada for a brother.
\end{quote}

\subsection*{You'll want to know: \kn{ಬಂಧು}}
\textbf{\kn{ಬಂಧು}} — \emph{bandhu} — \textquotedblleft{}kin; a relative\textquotedblright{}.

\begin{quote}
\kn{ಹಬ್ಬಕ್ಕೆ} \kn{ಎಲ್ಲಾ} \kn{ಬಂಧುಗಳು} \kn{ನಮ್ಮ} \kn{ಮನೆಗೆ} \kn{ಬಂದರು}.
\end{quote}

\begin{quote}
\emph{habbakke ellā bandhugaḷu namma manege bandaru.}
\end{quote}

\begin{quote}
All our relatives came to our house for the festival.
\end{quote}

Plural: \kn{ಬಂಧುಗಳು}, \emph{bandhugaḷu}.

\kn{ಬಂಧು}-\kn{ಬಳಗ}, relatives and friends.

\begin{cousinweb}
\textbf{\kn{ಬಂಧು}} (\emph{bandhu}) is borrowed from Sanskrit \emph{bandhu}, \textquotedblleft{}kinsman\textquotedblright{} (CDIAL 9144), from \emph{bandh}, \textquotedblleft{}to bind\textquotedblright{}.
\end{cousinweb}

\begin{grammarlens}[title={What you've built}]
One more word for describing what you lived and felt.
\end{grammarlens}

\subsection*{Your turn}
\begin{itemize}
\raggedright
  \item \emph{Say it:} \emph{bandhu}
  \item \emph{Say it:} \emph{bandhu}, once more
  \item \emph{Say it:} \emph{habbakke ellā bandhugaḷu namma manege bandaru.} — All our relatives came to our house for the festival.
  \item \emph{Recall:} say the Kannada for a cousin, then the Kannada for a brother, then say \emph{bandhu} again
\end{itemize}

\begin{tcolorbox}[breakable,colback=teal!4,colframe=teal!35!black,title={Before you move on}]
What does \kn{ಬಂಧು} mean? (\textquotedblleft{}Kin; a relative\textquotedblright{}.) Say it once more.
\end{tcolorbox}
\section[\texorpdfstring{\kn{ತವರು} (tavaru)}{ತವರು (tavaru)}]{\kn{ತವರು} (tavaru) — a woman's parental home}

\label{lesson:KA-C526-tavaru}



\begin{quote}
Before the new one: say the Kannada for a brother, then the Kannada for kin.
\end{quote}

\subsection*{You'll want to know: \kn{ತವರು}}
\textbf{\kn{ತವರು}} — \emph{tavaru} — \textquotedblleft{}a woman's parental home\textquotedblright{}.

\begin{quote}
\kn{ಹೆರಿಗೆಗೆ} \kn{ಅಕ್ಕ} \kn{ತವರಿಗೆ} \kn{ಬಂದಳು}.
\end{quote}

\begin{quote}
\emph{herigege akka tavarige bandaḷu.}
\end{quote}

\begin{quote}
My elder sister came to our parents' home for the birth.
\end{quote}

\kn{ತವರು} \kn{ಮನೆ}, the house a married woman was born in.

\begin{cousinweb}
\textbf{\kn{ತವರು}} (\emph{tavaru}) is usually traced to the native \emph{tavar}, \textquotedblleft{}those who are one's own people\textquotedblright{} (DEDR 3162); the path is not settled.
\end{cousinweb}

\begin{grammarlens}[title={What you've built}]
One more word for describing what you lived and felt.
\end{grammarlens}

\subsection*{Your turn}
\begin{itemize}
\raggedright
  \item \emph{Say it:} \emph{tavaru}
  \item \emph{Say it:} \emph{tavaru}, once more
  \item \emph{Say it:} \emph{herigege akka tavarige bandaḷu.} — My elder sister came to our parents' home for the birth.
  \item \emph{Recall:} say the Kannada for a brother, then the Kannada for kin, then say \emph{tavaru} again
\end{itemize}

\begin{tcolorbox}[breakable,colback=teal!4,colframe=teal!35!black,title={Before you move on}]
What does \kn{ತವರು} mean? (\textquotedblleft{}A woman's parental home\textquotedblright{}.) Say it once more.
\end{tcolorbox}
\section[\texorpdfstring{\kn{ಅವಳಿ} (avaḷi)}{ಅವಳಿ (avaḷi)}]{\kn{ಅವಳಿ} (avaḷi) — twins}

\label{lesson:KA-C526-avali}



\begin{quote}
Before the new one: say the Kannada for kin, then the Kannada for a woman's parental home.
\end{quote}

\subsection*{You'll want to know: \kn{ಅವಳಿ}}
\textbf{\kn{ಅವಳಿ}} — \emph{avaḷi} — \textquotedblleft{}twins\textquotedblright{}.

\begin{quote}
\kn{ನನ್ನ} \kn{ಅಕ್ಕನಿಗೆ} \kn{ಅವಳಿ} \kn{ಮಕ್ಕಳು} \kn{ಇದ್ದಾರೆ}.
\end{quote}

\begin{quote}
\emph{nanna akkanige avaḷi makkaḷu iddāre.}
\end{quote}

\begin{quote}
My elder sister has twins.
\end{quote}

\kn{ಅವಳಿ}-\kn{ಜವಳಿ}, a pair of twins.

\begin{cousinweb}
\textbf{\kn{ಅವಳಿ}} (\emph{avaḷi}) is the Kannada word for twins; its origin is not settled.
\end{cousinweb}

\begin{grammarlens}[title={What you've built}]
One more word for describing what you lived and felt.
\end{grammarlens}

\subsection*{Your turn}
\begin{itemize}
\raggedright
  \item \emph{Say it:} \emph{avaḷi}
  \item \emph{Say it:} \emph{avaḷi}, once more
  \item \emph{Say it:} \emph{nanna akkanige avaḷi makkaḷu iddāre.} — My elder sister has twins.
  \item \emph{Recall:} say the Kannada for kin, then the Kannada for a woman's parental home, then say \emph{avaḷi} again
\end{itemize}

\begin{tcolorbox}[breakable,colback=teal!4,colframe=teal!35!black,title={Before you move on}]
What does \kn{ಅವಳಿ} mean? (\textquotedblleft{}Twins\textquotedblright{}.) Say it once more.
\end{tcolorbox}
\section[\texorpdfstring{\kn{ಮರಿಮಗ} (marimaga)}{ಮರಿಮಗ (marimaga)}]{\kn{ಮರಿಮಗ} (marimaga) — a great-grandson}

\label{lesson:KA-C526-marimaga}



\begin{quote}
Before the new one: say the Kannada for a woman's parental home, then the Kannada for twins.
\end{quote}

\subsection*{You'll want to know: \kn{ಮರಿಮಗ}}
\textbf{\kn{ಮರಿಮಗ}} — \emph{marimaga} — \textquotedblleft{}a great-grandson\textquotedblright{}.

\begin{quote}
\kn{ಅಜ್ಜಿ} \kn{ತಮ್ಮ} \kn{ಮರಿಮಗನನ್ನು} \kn{ಎತ್ತಿಕೊಂಡು} \kn{ನಕ್ಕರು}.
\end{quote}

\begin{quote}
\emph{ajji tamma marimaganannu ettikoṇḍu nakkaru.}
\end{quote}

\begin{quote}
Grandmother picked up her great-grandson and laughed.
\end{quote}

\kn{ಮರಿ}, a young one + \kn{ಮಗ}, son. A great-granddaughter is \kn{ಮರಿಮಗಳು}.

\begin{cousinweb}
\textbf{\kn{ಮರಿಮಗ}} (\emph{marimaga}) is a Kannada compound built from \emph{mari}, \textquotedblleft{}a young one\textquotedblright{}, and \emph{maga}, \textquotedblleft{}son, male child\textquotedblright{} (DEDR 4616).
\end{cousinweb}

\begin{grammarlens}[title={What you've built}]
One more word for describing what you lived and felt.
\end{grammarlens}

\subsection*{Your turn}
\begin{itemize}
\raggedright
  \item \emph{Say it:} \emph{marimaga}
  \item \emph{Say it:} \emph{marimaga}, once more
  \item \emph{Say it:} \emph{ajji tamma marimaganannu ettikoṇḍu nakkaru.} — Grandmother picked up her great-grandson and laughed.
  \item \emph{Recall:} say the Kannada for a woman's parental home, then the Kannada for twins, then say \emph{marimaga} again
\end{itemize}

\begin{tcolorbox}[breakable,colback=teal!4,colframe=teal!35!black,title={Before you move on}]
What does \kn{ಮರಿಮಗ} mean? (\textquotedblleft{}A great-grandson\textquotedblright{}.) Say it once more.
\end{tcolorbox}
\section[\texorpdfstring{\kn{ಕಿರಿಯ} (kiriya)}{ಕಿರಿಯ (kiriya)}]{\kn{ಕಿರಿಯ} (kiriya) — younger; junior}

\label{lesson:KA-C526-kiriya}



\begin{quote}
Before the new one: say the Kannada for twins, then the Kannada for a great-grandson.
\end{quote}

\subsection*{You'll want to know: \kn{ಕಿರಿಯ}}
\textbf{\kn{ಕಿರಿಯ}} — \emph{kiriya} — \textquotedblleft{}younger; junior\textquotedblright{}.

\begin{quote}
\kn{ನನ್ನ} \kn{ಕಿರಿಯ} \kn{ಮಗ} \kn{ಶಾಲೆಗೆ} \kn{ಹೋಗುತ್ತಾನೆ}.
\end{quote}

\begin{quote}
\emph{nanna kiriya maga śālege hōguttāne.}
\end{quote}

\begin{quote}
My younger son goes to school.
\end{quote}

The opposite is \kn{ಹಿರಿಯ}, elder.

\begin{cousinweb}
\textbf{\kn{ಕಿರಿಯ}} (\emph{kiriya}) is native Dravidian: DEDR 1594 lists Kannada \emph{kiṟi}, \textquotedblleft{}small, little\textquotedblright{}, with \emph{kiṟiya}, \textquotedblleft{}younger\textquotedblright{}.
\end{cousinweb}

\begin{grammarlens}[title={What you've built}]
One more word for describing what you lived and felt.
\end{grammarlens}

\subsection*{Your turn}
\begin{itemize}
\raggedright
  \item \emph{Say it:} \emph{kiriya}
  \item \emph{Say it:} \emph{kiriya}, once more
  \item \emph{Say it:} \emph{nanna kiriya maga śālege hōguttāne.} — My younger son goes to school.
  \item \emph{Recall:} say the Kannada for twins, then the Kannada for a great-grandson, then say \emph{kiriya} again
\end{itemize}

\begin{tcolorbox}[breakable,colback=teal!4,colframe=teal!35!black,title={Before you move on}]
What does \kn{ಕಿರಿಯ} mean? (\textquotedblleft{}Younger; junior\textquotedblright{}.) Say it once more.
\end{tcolorbox}
